\documentclass[10pt,twocolumn]{article}

\usepackage[margin=0.72in,columnsep=0.25in]{geometry}
\usepackage[T1]{fontenc}
\usepackage[utf8]{inputenc}
\usepackage{microtype}
\usepackage{booktabs}
\usepackage{enumitem}
\usepackage[numbers,sort&compress]{natbib}
\usepackage[hidelinks]{hyperref}
\usepackage{xurl}

\setlength{\parindent}{1em}
\setlength{\parskip}{0pt}
\setlist[enumerate]{leftmargin=*,itemsep=2pt,topsep=3pt}

\title{\textbf{DEEDY: A Thai-Centric Multi-Agent Framework for Modeling\\
Emergent Public Sentiment Dynamics}}
\author{Sawit Koseeyaumporn, \\
\small Affiliation\\
\small \texttt{email@example.com}}
\date{}

\begin{document}
\maketitle

\begin{abstract}
Large language model (LLM) agents enable researchers to examine how local
interactions may produce population-level phenomena such as narrative
formation, polarization, and information diffusion. However, linguistically
plausible agent behavior does not by itself establish that a simulation is a
valid proxy for human behavior, particularly in a culturally and
linguistically specific setting. This paper proposes DEEDY---Digital
Environment for Emergent Dynamics and Yield---a Thai-centric multi-agent
framework for modeling and evaluating emergent public sentiment dynamics.
DEEDY extends the MiroFish/OASIS pipeline with Thai-language data processing,
culturally grounded actor construction, and validation against observed online
discussions. The framework transforms news and social-trend data into a
graph-grounded environment, generates heterogeneous actors, executes
multi-round social interactions, and produces auditable simulation traces and
reports. For each event, only scenario-construction data are exposed to the
simulation, while real posts and comments are held out as ground truth.
Simulated and observed discourse are compared at the content, population, and
diffusion levels using sentiment and stance distributions, narrative
correspondence, temporal propagation, interaction structure, and Thai
linguistic and cultural appropriateness. DEEDY is therefore positioned as an
empirically testable instrument for exploratory social simulation rather than
as an unquestioned substitute for Thai participants or a point predictor of
public opinion.
\end{abstract}

\noindent\textbf{Keywords:} LLM agents, agent-based modeling, social
simulation, public sentiment, Thai NLP, information diffusion

\section{Introduction}
\label{sec:introduction}

Public sentiment on social media is not merely the sum of independent
opinions. It develops through repeated exposure, interpersonal influence,
platform recommendation, and feedback between individual actions and the
information environment. Consequently, phenomena such as narrative formation,
polarization, herding, and misinformation diffusion are emergent: they arise
from many local interactions and may not be recoverable by analyzing isolated
messages alone. Agent-based modeling (ABM) offers a bottom-up approach to this
problem by representing heterogeneous actors, their environment, and the rules
through which interactions accumulate into population-level outcomes
\citep{macal2010tutorial,gao2024survey}.

Traditional ABMs usually encode behavior through explicit rules, transition
probabilities, or utility functions. These mechanisms are interpretable but
can be difficult to specify for open-ended communication, where action depends
on long conversational histories, social roles, and contextual meaning. LLM
agents provide a complementary mechanism: an agent can condition its actions
on a persona, memory, retrieved evidence, and the current social environment.
Early systems showed that such agents can generate believable routines and
coordinated group behavior \citep{park2023generative}, while social-media
prototypes demonstrated the generation of posts, replies, and moderation
interactions at community scale \citep{park2022social}. These capabilities
make LLM agents attractive for studying discourse dynamics that are difficult
to express as fixed behavioral rules.

Scalable social simulation systems have subsequently broadened this approach.
$S^3$ modeled emotion, attitude, and interaction behavior using real social
network data \citep{gao2023s3}. OASIS introduced a dynamic social-media
environment with configurable recommendation mechanisms, diverse actions, and
support for simulations of up to one million agents
\citep{yang2024oasis}. GenSim emphasized reusable scenario functions,
large-scale execution, and error correction \citep{tang2025gensim}, while
AgentSociety coupled generative agents with social, urban, and economic
environments \citep{piao2025agentsociety}. MiroFish builds on OASIS as an
open-source scenario pipeline that extracts entities and relationships from
``reality seed'' documents, constructs graph-based memory, generates agent
personas, runs social simulations, and synthesizes a report
\citep{mirofish2026}.

Despite this progress, two gaps limit the direct use of these systems for Thai
public-sentiment research. First, localization is not equivalent to translating
prompts into Thai. Thai is commonly written without explicit word boundaries,
and social-media text adds slang, non-standard spelling, emoji, transliterated
loanwords, and Thai--English code-mixing. Thai NLP remains comparatively
under-resourced, especially for downstream semantic tasks
\citep{arreerard2022survey,phatthiyaphaibun2023pythainlp}. Thai-specific
models such as WangchanBERTa and Typhoon demonstrate that language-specific
pretraining, tokenization, and instruction tuning improve Thai processing
\citep{lowphansirikul2021wangchanberta,pipatanakul2023typhoon}. Nevertheless,
recent benchmarking also finds that performance can decline substantially for
Northern, Northeastern, and Southern Thai varieties compared with Standard
Thai \citep{limkonchotiwat2025dialects}. A simulation that ignores these
features may produce fluent Standard Thai while misrepresenting the linguistic
and social heterogeneity of the population of interest.

Second, believable interaction is not evidence of behavioral validity. LLMs
can sometimes reproduce subgroup-level response distributions when carefully
conditioned \citep{argyle2023out}, and interview-grounded agents have shown
promising agreement with the survey responses of the individuals they
represent \citep{park2024thousand}. However, other studies report diminished
diversity and sensitivity to prompting or answer order
\citep{park2024diversity}. A critical review of generative ABMs concludes that
many studies still rely on face validity or outcome measures that are only
loosely connected to the mechanisms being modeled
\citep{larooij2026validation}. Thus, a Thai-language simulation should not be
treated as a proxy for Thai society unless its outputs are compared with
held-out observations at both micro and macro levels.

To address these gaps, we propose \textbf{DEEDY} (\textbf{D}igital
\textbf{E}nvironment for \textbf{E}mergent \textbf{D}ynamics and
\textbf{Y}ield), a Thai-centric extension of the MiroFish/OASIS workflow.
DEEDY separates \emph{scenario-construction evidence} from
\emph{validation evidence}. News reports and pre-event contextual material are
used to construct the knowledge graph, actors, and environment; the
corresponding real social posts, comments, and temporal interaction records are
withheld from the agents and used only after simulation. This design reduces
direct answer leakage and makes it possible to ask whether emergent simulation
outcomes correspond to observed Thai discourse.

DEEDY evaluates correspondence across multiple levels rather than relying on a
single LLM-as-a-judge score. Content-level evaluation measures sentiment,
stance, topic, narrative, and linguistic properties. Population-level
evaluation compares distributions, subgroup differences, polarization, and
uncertainty across repeated runs. Dynamic evaluation compares diffusion scale,
depth, breadth, timing, and interaction structure, following the broader
principle of validating generative ABMs against real ground truth
\citep{gao2024survey,yang2024oasis}. Thai human assessment and LLM-assisted
judging complement, but do not replace, quantitative comparison. The intended
claim is therefore bounded: DEEDY is designed to reveal when a Thai-centric
LLM-agent simulation reproduces observed patterns, when it fails, and how
those conclusions change under alternative models and simulation settings.

\section{Research Objectives and Contributions}
\label{sec:objectives}

This study has three contribution-level objectives:

\begin{enumerate}
    \item \textbf{Develop a Thai-centric multi-agent framework.}
    Extend the MiroFish/OASIS pipeline with Thai-aware data processing,
    graph-grounded actor and environment construction, and representations of
    linguistic and sociocultural heterogeneity, including informal language,
    code-mixing, and regional varieties where supported by the data.

    \item \textbf{Model emergent Thai public-sentiment dynamics.}
    Use interactions among heterogeneous LLM agents to examine how individual
    exposure and response may aggregate into sentiment change, stance
    formation, narrative competition, polarization, and information diffusion
    under different social-network and recommendation conditions.

    \item \textbf{Establish an empirical validation framework.}
    Evaluate simulated outcomes against held-out Thai posts, comments, and
    interaction traces using content-, population-, and diffusion-level
    metrics, repeated runs, uncertainty reporting, and Thai human evaluation,
    thereby identifying both the framework's useful operating conditions and
    its limitations as a proxy for human behavior.
\end{enumerate}

\section{Background}
\label{sec:background}

\subsection{Agent-Based and Generative Social Simulation}

An ABM specifies a population of autonomous units whose states and actions
change through interaction with one another and with an environment. The
researcher observes whether macro-level regularities emerge from these
micro-level mechanisms \citep{macal2010tutorial}. In a social-media setting,
an agent may read recommended content, create a post, reply, follow another
user, or remain inactive. The platform environment stores content and
relationships, advances time, and determines which information each agent can
observe. Recommendation policies are part of the model because they mediate
exposure and can therefore affect diffusion and polarization.

LLM-based ABMs replace or augment some hand-written decision rules with
language-model inference. The usual architecture combines a relatively stable
profile, a changing state, memory, perception of available content, and an
action-selection prompt. This improves expressive capacity but introduces new
sources of uncertainty. Behavior can change with the base model, prompt
wording, decoding parameters, memory design, or action ordering. Furthermore,
agents produced by a shared model may be less heterogeneous than their persona
descriptions suggest \citep{park2024diversity}. Repeated simulation,
sensitivity analysis, and explicit reporting of model configurations are
therefore part of the scientific method rather than implementation details.

\subsection{The OASIS and MiroFish Foundations}

OASIS is the simulation substrate most directly relevant to DEEDY. Its
environment maintains user profiles, posts, and social relationships; its
information filter and recommendation modules determine visible content; and
its agents select among social actions over discrete time steps
\citep{yang2024oasis}. OASIS evaluated information propagation using scale,
depth, maximum breadth, and normalized root-mean-square error against real
propagation curves. It also studied polarization and herding, illustrating the
importance of combining behavioral experiments with real-data comparison.

MiroFish adds a document-to-simulation workflow on top of this substrate. Its
official implementation describes four stages: graph building from seed
material, environment and persona construction, multi-agent simulation, and
report generation \citep{mirofish2026}. This workflow is useful for rapid
scenario construction, but a generated report is an interpretation of the
simulated world, not an independent validation of it. DEEDY therefore retains
the graph-grounded construction and social-simulation components while adding
a formal separation between inputs and held-out ground truth, Thai-specific
representations, repeated-run analysis, and explicit validation artifacts.

\subsection{Thai Language and Online Discourse}

Thai presents both computational and sociolinguistic challenges. The absence
of explicit word delimiters makes segmentation a substantive modeling choice,
not a neutral preprocessing step \citep{chormai2020segmentation}. Social-media
messages further contain creative spelling, repeated characters, emoji,
hashtags, abbreviations, and mixed scripts. Code-switching can also serve
pragmatic functions such as expressing identity, politeness, group membership,
or emphasis \citep{kongkerd2016codeswitching}; consequently, normalizing every
message into formal monolingual Thai may erase socially meaningful signals.

The Thai NLP ecosystem nevertheless provides important foundations.
PyThaiNLP offers open tools and resources for tokenization and other Thai NLP
tasks \citep{phatthiyaphaibun2023pythainlp}. WangchanBERTa showed the value of
Thai-specific preprocessing and pretraining on data that include news and
social media \citep{lowphansirikul2021wangchanberta}; Typhoon extended this
direction to instruction-following Thai LLMs
\citep{pipatanakul2023typhoon}. The Wisesight Sentiment Corpus provides
26,737 informal Thai messages labeled as positive, neutral, negative, or
question \citep{wisesight2020}. These resources can support preprocessing,
sentiment measurement, and baseline evaluation, but they do not by themselves
model longitudinal interaction or guarantee representativeness of Thai public
opinion. DEEDY consequently treats language resources as measurement
instruments and calibration aids, while constructing event-specific validation
sets from observed discussions under documented sampling and privacy
procedures.

\section{Related Work}
\label{sec:related}

\subsection{LLM Agents for Individual and Social Behavior}

Generative Agents introduced a memory--reflection--planning architecture that
produced believable individual routines and emergent coordination in a small
interactive world \citep{park2023generative}. Social Simulacra applied related
capabilities to populated prototypes of online communities, generating
members, posts, replies, and antisocial behavior in response to a system
description \citep{park2022social}. These systems established the expressive
potential of LLM-mediated simulation but primarily evaluated believability and
design utility.

Research on simulated human samples provides a complementary line of evidence.
Argyle et al.\ introduced \emph{algorithmic fidelity}, comparing
demographically conditioned model responses with human subgroup distributions
\citep{argyle2023out}. Park et al.\ constructed agents from qualitative
interviews with 1,052 individuals and evaluated them using survey,
personality, and experimental measures \citep{park2024thousand}. These studies
show that empirical grounding can improve correspondence, but they primarily
concern individual response replication rather than the evolution of public
discourse through networked interaction.

\subsection{Large-Scale Social-Media Simulation}

$S^3$ was among the early systems to connect LLM-based perception and action
with emotion, attitude, and social-network interaction
\citep{gao2023s3}. OASIS expanded scalability and modeled platform mechanisms,
including dynamic networks and recommendation systems
\citep{yang2024oasis}. GenSim provided abstractions for diverse scenarios and
mechanisms for correcting simulation errors
\citep{tang2025gensim}, whereas AgentSociety integrated broader social, urban,
and economic spaces \citep{piao2025agentsociety}. DEEDY does not compete with
these systems primarily on agent count. Its intended contribution is a
language- and culture-specific construction-and-validation layer for Thai
public discourse, implemented as an extension of the MiroFish/OASIS workflow.

\subsection{Thai Sentiment and Language Modeling}

Prior Thai NLP research has concentrated on resources, tokenization, text
classification, and language-model adaptation. Surveys describe strong
progress in upstream processing but fewer resources for downstream semantic
analysis \citep{arreerard2022survey}. Thai social-text models demonstrate that
domain-specific pretraining and preprocessing improve classification
\citep{horsuwan2019socialtext}, and WangchanBERTa provides a widely used
Thai-specific encoder \citep{lowphansirikul2021wangchanberta}. More recent Thai
LLMs improve generation and instruction following
\citep{pipatanakul2023typhoon}, yet dialect benchmarks indicate persistent
performance disparities \citep{limkonchotiwat2025dialects}.

Most of this work evaluates isolated texts or task responses. Sentiment
classification can label a message, but it does not explain how exposure,
reply structure, social influence, or recommendation produce a population
trajectory. DEEDY connects Thai NLP measurement with an interactive
multi-agent environment and evaluates both the generated language and the
resulting collective dynamics.

\subsection{Validation of Generative Social Simulation}

Validation remains the central methodological distinction between a compelling
demo and a scientific simulation. OASIS compared simulated and real
propagation curves and reported structural limitations in reproducing
diffusion depth \citep{yang2024oasis}. Broader surveys recommend validation at
both individual and population levels using real ground truth
\citep{gao2024survey}. Conversely, evidence of low response diversity and
order sensitivity shows why surface realism or a single successful run is
insufficient \citep{park2024diversity}. The systematic review by Larooij and
T\"ornberg further argues that black-box mechanisms, cultural bias, and
stochasticity can make generative ABMs harder rather than easier to validate
\citep{larooij2026validation}.

DEEDY operationalizes these lessons through event-level holdout, repeated
runs, quantitative comparisons, human assessment, and sensitivity analysis.
Its evaluation is designed to distinguish at least three questions: whether
individual outputs are linguistically and culturally plausible; whether
aggregate sentiment, stance, and narrative distributions correspond to the
observed population; and whether the temporal and network processes that
produce those outcomes resemble real diffusion. This separation prevents a
high score on textual naturalness from being misinterpreted as evidence that
the simulated social process is valid.

\section{Scope Statement}
\label{sec:scope}

DEEDY is intended for exploratory and comparative analysis of Thai online
discourse. Unless supported by prospective validation across multiple events,
its outputs should not be described as forecasts of how Thai people will
behave. The framework models the sampled online environment, not the entire
Thai population, and any empirical study must document data provenance,
platform coverage, demographic limitations, de-identification, and compliance
with applicable research-ethics and platform requirements.

\bibliographystyle{unsrtnat}
\bibliography{deedy_references}

\end{document}
